% ============================================================
% DISCUSSION — target: 1000--1500 words
% ============================================================

\section{Discussion}

% Paragraph 1: Summary and takeaway
The primary aim of this study was to determine whether the low CoT faithfulness
rates observed in proprietary reasoning models generalize across the open-weight
reasoning model ecosystem. [TODO: State primary finding. ``Our results demonstrate
that...'' State overall takeaway.]

% Paragraph 2: Cross-family variation (H1)
[TODO: Restate H1 finding. Interpret: what explains the cross-family differences?
Relate to training methodology differences (GRPO vs.\ DPO vs.\ distillation vs.\
data-centric). Compare with Chen et al.'s finding that DeepSeek-R1 (39\%) outperforms
Claude 3.7 Sonnet (25\%). To our knowledge, this is the first study to demonstrate
systematic differences in CoT faithfulness across model families.]

% Paragraph 3: Hint-type patterns (H2)
[TODO: Restate H2 finding. Interpret: why might certain hint types be systematically
underreported? Relate to Turpin et al.'s categories of biasing features. Discuss
implications for monitoring---if metadata and grader hints are consistently
unfaithful, these represent the highest-risk categories for safety monitoring.]

% Paragraph 4: Scaling
[TODO: Interpret scaling results. Does larger = more faithful? Relate to the general
finding that larger models are more capable but not necessarily more aligned. Compare
with the 8B distilled model vs.\ full DeepSeek-R1.]

% Paragraph 5: CoT length
[TODO: Interpret CoT length findings. If replicated, unfaithful CoTs being longer
suggests active rationalization rather than simple omission---models may construct
elaborate alternative justifications when they don't want to reveal the true influence.
Relate to Chen et al.'s ``obfuscation'' interpretation.]

% Limitations
\subsection*{Limitations}

Several limitations should be considered when interpreting these results.

First, our faithfulness measurement is inherently conservative: we only identify a
CoT as faithful if it \emph{explicitly} mentions the hint, but models may
acknowledge hint influence through indirect or paraphrased language that our
classifier misses. Our hand-labeled validation set provides a bound on this error,
but some faithful responses may be misclassified as unfaithful.

Second, we evaluate using API-served models, which may differ from locally-hosted
versions due to provider-specific quantization, system prompts, or inference
optimizations. While we control for generation parameters
($\text{temperature}=0.0$, $\text{seed}=42$), exact reproducibility across
API providers cannot be guaranteed.

Third, our evaluation covers multiple-choice questions from MMLU and GPQA only.
Faithfulness behavior may differ on open-ended generation tasks, code
generation, or mathematical reasoning where the answer format is unconstrained.

Fourth, we test 500 questions per model, which may provide insufficient statistical
power for fine-grained comparisons within specific MMLU subject categories or for
rare hint types with low influence rates.

[TODO: Address additional limitations discovered during experiments.]

% Integration: implications
\subsection*{Implications for CoT Safety Monitoring}

[TODO: Theoretical implications---what does cross-family faithfulness variation
tell us about the relationship between training methodology and reasoning
transparency? Practical implications---which models are most suitable for
deployments that rely on CoT monitoring? How should safety teams adjust their
monitoring strategies given these findings?

Relate to Baker et al.'s argument that CoT can be informative despite
unfaithfulness. Our cross-family data may help quantify \emph{how} informative
CoT monitoring is for different model families.

Discuss hybrid behavioral + mechanistic monitoring as a future direction,
following the argument that CoT visibility alone is insufficient for reliable
safety monitoring.]

% Future directions
\subsection*{Future Directions}

[TODO: (1)~Extend to open-ended generation tasks beyond multiple choice.
(2)~Combine CoT analysis with probing-based approaches to measure the gap
between surface-level and internal faithfulness. (3)~Investigate whether
faithfulness can be improved through training-time interventions (e.g.,
faithfulness-aware RLHF) without sacrificing task performance. (4)~Longitudinal
tracking of faithfulness as models are updated across versions.]
